\section{Related work}
\label{sec:related}
Diagonal scaling has a long history, including the logarithmic method of Curtis and
Reid \cite{curtis1972scaling}. Elble and Sahinidis \cite{elble2012scaling} compare
scalers for simplex and find that additional complexity need not improve performance.
Klotz \cite{klotz2014numerics} relates numerical instability to formulation and finite
precision. Optimal diagonal preconditioning is a broader problem: Qu et al.
\cite{qu2024diagonal} use quasiconvex and semidefinite formulations, while Demmel
\cite{demmel2023block} gives block-Jacobi guarantees. Our triangular calculation is
an explanatory scalar special case, not a general preconditioning algorithm.

Robustness is also an established objective in MIP scaling. Berthold and Hendel
\cite{berthold2021scaling} learn to select between standard and Curtis--Reid scaling
in Xpress from an indicator of numerical difficulties. Their implementation uses
power-of-two factors and their evaluation examines inconsistencies across seeds.
Our external construction instead states row-level admissibility conditions; it does
not predict or replace internal solver scaling.

Modeling environments already support generator-level information: Pyomo's
\texttt{ScaleModel} applies component factors \cite{pyomoScaleModel}, GAMS has
variable and equation scale attributes \cite{gamsScaling}, and IDAES derives and
propagates scales across submodels \cite{idaesScaling,idaesCustomScaler}. Plasmo.jl
retains local problems and linking constraints \cite{jalving2022plasmo}. Consequently,
using modeling metadata before export is not new. Our identifiability result concerns
the loss of one specific block statistic after a preceding normalization.

Tolerance-based scaling also has direct precedent. Gurobi documents absolute
feasibility and integrality tolerances and the treatment of very small coefficients
\cite{gurobiScaling}, and recommends external scaling that makes tolerance-sized
violations negligible in application units \cite{gurobiExternalScaling}. Pritchett,
Howell and Grebow \cite{pritchett2017scaling} similarly scale constraints using
physical accuracy requirements and SNOPT's tolerance. Section~\ref{sec:contract}
combines this known transfer with coefficient limits, identifies incompatible
specifications, and solves the constrained centering problem. Integer rounding and
exact binary64 export impose additional restrictions.

Finally, Cheung, Gleixner and Steffy \cite{cheung2017certificates} provide certificates
for integer-programming results, and Hoen and Gleixner \cite{hoen2024correctness}
analyze individual branch-and-bound decisions. Our rational residual checks concern
primal acceptance; independent enumeration establishes optimality only on the small
binary instances. The complementary attribution audit checks actual exports under
matched kernels and coverage. These contributions concern identifiable mechanisms
and explicit numerical contracts, without claiming universal performance superiority.
